\documentclass[../DoAn.tex]{subfiles}
\begin{document}

Chương này trình bày quá trình khảo sát hiện trạng và phân tích yêu cầu của hệ thống IoT giám sát người cao tuổi. Phần \ref{section:2.1} khảo sát hiện trạng bài toán phát hiện té ngã, bao gồm các sản phẩm thương mại đã triển khai thực tế (\ref{subsection:2.1.1}) và tổng quan các hướng nghiên cứu học thuật theo ba nhóm phương pháp chính (\ref{subsection:2.1.2}). Phần \ref{section:2.2} và \ref{section:2.3} lần lượt trình bày tổng quan chức năng và đặc tả chi tiết các ca sử dụng cốt lõi. Phần \ref{section:2.4} tổng hợp các yêu cầu phi chức năng về độ trễ, bảo mật và khả năng mở rộng.

\section{Khảo sát hiện trạng}
\label{section:2.1}

Bài toán phát hiện té ngã tự động bằng thiết bị đeo wearable đã được nghiên cứu rộng rãi trong hơn hai thập kỷ qua do tầm quan trọng y tế cấp thiết của nó. Tổ chức Y tế Thế giới (WHO) ước tính khoảng 684.000 ca tử vong mỗi năm liên quan đến té ngã, trong đó người cao tuổi chiếm tỷ lệ cao nhất \cite{mekruksavanich2024resnet}. Điều đặc biệt nghiêm trọng là khoảng 50\% người cao tuổi nằm bất động dưới sàn sau khi ngã quá một giờ sẽ tử vong trong vòng sáu tháng, bất kể chấn thương trực tiếp có nghiêm trọng hay không \cite{mekruksavanich2024resnet}. Điều này đặt ra yêu cầu bắt buộc về việc rút ngắn thời gian phát hiện và can thiệp y tế.

Về mặt kỹ thuật, bài toán phát hiện té ngã được phân chia thành hai pha xử lý chính. \textbf{Phát hiện trước va chạm} (Pre-Impact Fall Detection --- PIFD) nhận diện chuyển động ngã trong khoảng thời gian từ khi bắt đầu mất thăng bằng đến khi cơ thể chạm đất, thường chỉ kéo dài từ 100 đến 500 mili giây. Pha này đòi hỏi suy luận hoàn toàn trên thiết bị (on-device inference) với độ trễ cực thấp, không thể dựa vào kết nối đám mây. \textbf{Phát hiện sau va chạm} (Post-Impact Fall Detection) xác nhận cú ngã đã xảy ra sau khi cơ thể tiếp đất, thường kết hợp với cơ chế chờ phản hồi của người dùng trước khi phát cảnh báo. Pha này có độ trễ cho phép lớn hơn song yêu cầu kết nối mạng để thông báo đến người giám sát. Đồ án này tập trung vào pha thứ hai: phát hiện sau va chạm kết hợp với truyền thông cảnh báo qua 4G LTE đến dashboard y tế.

\subsection{Hiện trạng các sản phẩm thương mại}
\label{subsection:2.1.1}

Sự cấp thiết của bài toán đã thúc đẩy sự ra đời của nhiều sản phẩm thương mại, tập trung vào hai hướng tiếp cận tương ứng với hai pha phát hiện trên.

\textbf{Nhóm thiết bị túi khí bảo vệ chủ động (PIFD).} Hướng tiếp cận này tích hợp IMU on-device với cơ chế bơm túi khí khí nén để bảo vệ các vùng xương dễ gãy trong khoảng thời gian trước khi cơ thể chạm đất. Hövding\footnote{\url{https://hovding.com}} (Hövding Sverige AB, Thụy Điển) là vòng đeo cổ dành cho người đi xe đạp, tích hợp túi khí bao phủ toàn bộ đầu và cổ; hệ thống xử lý dữ liệu IMU ở tần số 200~Hz ngay trên vi điều khiển nhúng và kích hoạt túi khí trong vòng 0,1 giây --- nhanh hơn phản xạ bảo vệ tự nhiên của con người. Trong lĩnh vực thể thao mô-tô, Alpinestars Tech-Air\footnote{\url{https://www.alpinestars.com/pages/tech-air}} (Alpinestars, Ý) là hệ thống áo giáp túi khí tích hợp cho người lái xe máy, phát hiện va chạm và bơm phồng toàn bộ vùng vai, ngực và lưng trong khoảng 25~mili giây, sử dụng thuật toán phân tích IMU nhiều trục chạy hoàn toàn độc lập với điện thoại. Chuyển sang lĩnh vực chăm sóc người cao tuổi, Helite\footnote{\url{https://www.helite.com}} (Pháp) cung cấp hai sản phẩm có định vị khác nhau: \textbf{Hip'Air} là đai hông túi khí điện tử, bơm phồng hai túi khí bảo vệ khớp hông trong khoảng 80~mili giây sau khi phát hiện ngã; \textbf{HipSafe} là phiên bản thụ động hơn, sử dụng tấm đệm hấp thụ lực cơ học tích hợp sẵn trong đai, không cần nguồn điện hay thuật toán, phù hợp với người cao tuổi không quen dùng thiết bị điện tử. TangoBelt\footnote{\url{https://www.tangobelt.com}} là đai bụng chuyên dụng cho người cao tuổi, kết hợp cả hai cơ chế trong cùng một nền tảng phần cứng: phiên bản túi khí bơm phồng bảo vệ hông khi phát hiện ngã (tương tự Hip'Air), trong khi phiên bản kết nối gửi cảnh báo tức thì đến người thân hoặc nhân viên chăm sóc qua Bluetooth Low Energy (BLE) và ứng dụng di động đi kèm. Điểm nổi bật là khả năng tích hợp cả bảo vệ vật lý lẫn luồng thông báo xã hội (social alerting) trên cùng một thiết bị. Hạn chế của phiên bản kết nối là phụ thuộc vào điện thoại trung gian và vùng phủ BLE (khoảng 10 mét), chưa đáp ứng được môi trường cơ sở điều dưỡng quy mô lớn với nhiều bệnh nhân phân tán trên nhiều phòng.

\textbf{Nhóm đồng hồ thông minh (Post-Impact FD).} Apple Watch Series~4 trở lên\footnote{\url{https://support.apple.com/en-us/111839}} tích hợp tính năng Fall Detection sử dụng gia tốc kế và con quay hồi chuyển kết hợp với mô hình phân loại học máy chạy hoàn toàn trên thiết bị (on-device). Khi phát hiện cú ngã mạnh, đồng hồ hiển thị thông báo yêu cầu xác nhận; nếu người dùng không phản hồi trong vòng 60 giây, hệ thống tự động gọi cấp cứu và gửi vị trí GPS đến danh sách liên lạc khẩn cấp. Samsung Galaxy Watch cung cấp tính năng tương đương. Nhóm sản phẩm này phù hợp cho người dùng cá nhân di động nhưng không đáp ứng nhu cầu của cơ sở y tế vì thiếu dashboard giám sát tập trung đa bệnh nhân, không lưu trữ lịch sử vận động có cấu trúc và không tích hợp quy trình xác nhận xử lý sự cố theo chuẩn nghiệp vụ y tế.

% =====================================================================
% HÌNH MINH HỌA SẢN PHẨM THƯƠNG MẠI (lưới 2x2)
% Bỏ comment toàn bộ khối \begin{figure}...\end{figure} bên dưới
% sau khi đặt 4 file ảnh vào thư mục Hinhve/:
%   - hovding.jpg      : vòng đeo cổ Hövding đang bơm phồng túi khí
%   - alpinestars.jpg  : áo giáp Alpinestars Tech-Air (mô-tô)
%   - hipair.jpg       : đai hông Hip'Air với túi khí hai bên
%   - tangobelt.jpg    : thiết bị đai bụng TangoBelt kèm app di động
% Nguồn ảnh gợi ý: trang press/media của từng hãng hoặc ảnh từ paper
% học thuật trích dẫn sản phẩm (thường miễn phí sử dụng trong luận văn).
% =====================================================================
\begin{figure}[H]
    \centering
    % --- Hàng trên: nhóm túi khí thể thao ---
    \begin{subfigure}[b]{0.46\textwidth}
        \centering
        \includegraphics[width=\textwidth]{Hinhve/hovding.jpg}
        \caption{Hövding 3 --- vòng cổ túi khí cho người đi xe đạp
                 (Nguồn: hovding.com)}
        \label{fig:hovding}
    \end{subfigure}
    \hfill
    \begin{subfigure}[b]{0.46\textwidth}
        \centering
        \includegraphics[width=\textwidth]{Hinhve/alpinestars.jpg}
        \caption{Alpinestars Tech-Air --- áo giáp túi khí dành cho người lái mô-tô
                 (Nguồn: alpinestars.com)}
        \label{fig:alpinestars}
    \end{subfigure}
    \vspace{0.5em}
    % --- Hàng dưới: nhóm chăm sóc người cao tuổi ---
    \begin{subfigure}[b]{0.46\textwidth}
        \centering
        \includegraphics[width=\textwidth]{Hinhve/hipair.jpg}
        \caption{Hip'Air (Helite) --- đai hông túi khí bảo vệ khớp hông
                 người cao tuổi (Nguồn: helite.com)}
        \label{fig:hipair}
    \end{subfigure}
    \hfill
    \begin{subfigure}[b]{0.46\textwidth}
        \centering
        \includegraphics[width=\textwidth]{Hinhve/tangobelt.jpg}
        \caption{TangoBelt --- đai bụng IoT kết nối BLE cảnh báo té ngã
                 cho người cao tuổi (Nguồn: tangobelt.com)}
        \label{fig:tangobelt}
    \end{subfigure}
    \caption{Các sản phẩm thiết bị đeo phát hiện té ngã thương mại tiêu biểu,
             phân loại theo đối tượng sử dụng và cơ chế bảo vệ.}
    \label{fig:commercial_products}
\end{figure}

Khoảng trống rõ ràng từ phân tích trên là chưa có giải pháp nào kết hợp đồng thời: suy luận TinyML on-device, truyền thông 4G LTE độc lập Wi-Fi, dashboard y tế thời gian thực đa bệnh nhân và lưu trữ lịch sử vận động chuỗi thời gian --- đây chính là mục tiêu mà đồ án này hướng đến.

\subsection{Các phương pháp nghiên cứu giải quyết bài toán}
\label{subsection:2.1.2}

Song song với sự phát triển của các sản phẩm thương mại, cộng đồng học thuật đã đề xuất nhiều hướng tiếp cận khác nhau để giải quyết bài toán phát hiện té ngã bằng cảm biến IMU. Hệ thống trong đồ án này kế thừa và vận dụng tri thức từ ba dòng nghiên cứu chính được khảo sát dưới đây.

\subsubsection{Phương pháp dựa trên ngưỡng}
\label{subsubsection:2.1.2.a}

Hướng tiếp cận dựa trên ngưỡng (threshold-based) là nhóm phương pháp lâu đời nhất và có chi phí tính toán thấp nhất, phù hợp với vi điều khiển năng lượng thấp. Nguyên lý chung là so sánh các đặc trưng tín hiệu IMU (biên độ gia tốc, căn bậc hai trung bình --- RMS, góc nghiêng) với một tập ngưỡng được định nghĩa trước; khi tín hiệu vượt ngưỡng, hệ thống kết luận có sự kiện té ngã.

\begin{enumerate}[label=\arabic*)]

\item Alsaadi và cộng sự \cite{alsaadi2024logical} đề xuất hệ thống nhận diện hoạt động con người kết hợp suy luận logic dựa trên đặc trưng vật lý của chuyển động với cơ chế hiệu chỉnh ngưỡng thích ứng bằng hồi quy véc-tơ hỗ trợ (Support Vector Regression --- SVR). Dữ liệu được thu thập từ mạng lưới bốn cảm biến IMU 9 trục đặt tại đầu, cổ tay, thắt lưng và bàn chân của 20 tình nguyện viên, đồng bộ hóa ở tần số 50 Hz. Thay vì dùng mô hình hộp đen, hệ thống phân loại 12 hoạt động hằng ngày thông qua một ma trận quan hệ hành động và cây logic dựa trên giá trị RMS, biên độ, và chu kỳ tín hiệu. Điểm sáng tạo nổi bật là cơ chế thích ứng cá nhân hóa: chỉ cần người dùng mới đi bộ thử hai bước, thuật toán SVR sẽ nội suy toàn bộ tập ngưỡng cho 11 hành vi còn lại, giải quyết trực tiếp bài toán khái quát hóa chéo người dùng (cross-person generalization) mà không cần thu thập dữ liệu té ngã từ đối tượng mới. Hệ thống đạt độ chính xác trung bình 90,8--92,1\% trong kịch bản chéo người dùng, vượt trội SVM (84,7\%) và DT (80,2\%). Đối với đồ án, pipeline tiền xử lý "bộ lọc bù bậc nhất + RMS triggering" là công thức có thể triển khai trực tiếp trên vi điều khiển ESP32; hạn chế chính là yêu cầu đồng thời bốn thiết bị đeo, không phù hợp với kịch bản một cảm biến duy nhất.

\item Tseng và cộng sự \cite{tseng2025wearable} phát triển thiết bị đeo tích hợp vi điều khiển, cảm biến IMU, GPS và NB-IoT (Narrowband Internet of Things) nhằm khắc phục giới hạn nhận thức vị trí của các hệ thống phát hiện té ngã trong nhà. Thuật toán được thiết kế theo mô hình máy trạng thái hữu hạn (Finite-State Machine --- FSM), có khả năng lọc tám loại hoạt động hằng ngày và phân biệt bốn hướng ngã (trước, sau, trái, phải). Khi phát hiện té ngã, thiết bị lập tức gửi cảnh báo kèm tọa độ GPS qua NB-IoT đến máy chủ, giải quyết bài toán cứu hộ trong môi trường ngoài trời. Thử nghiệm thu thập 6.750 mẫu từ 15 người tham gia, hệ thống đạt độ nhạy 97,9\%, độ đặc hiệu 99,9\% và độ chính xác tổng thể 99,7\%. Nghiên cứu khẳng định tính khả thi của FSM trên phần cứng nhúng năng lượng thấp và kiến trúc truyền thông hướng sự kiện (event-driven) --- hướng tiếp cận này được đồ án vận dụng làm cơ sở thiết kế FSM phát hiện ngã trên firmware ESP32-S3.

\item Booranawong và cộng sự \cite{booranawong2025stairs} giới thiệu hệ thống giám sát và phân loại hoạt động trên cầu thang theo thời gian thực, sử dụng module không dây XBee3 chuẩn IEEE 802.15.4 ở băng tần 2,4 GHz kết hợp gia tốc kế ba trục GY-521 đeo tại thắt lưng. Pipeline xử lý gồm tính véc-tơ độ lớn tín hiệu (SVM), lọc nhiễu bằng bộ lọc EWMA ($\alpha = 0{,}3$), trích xuất đặc trưng trong cửa sổ 15 mẫu (1,5 giây ở 10 Hz), và phân loại bằng K lân cận gần nhất (KNN). Kết quả đáng chú ý là hệ thống nhận diện cú ngã với độ chính xác 100\%, đồng thời đạt tỷ lệ gói tin thành công (PDR) 100\% cả trong điều kiện có và không có tầm nhìn thẳng. Nghiên cứu chứng minh rằng chỉ một đặc trưng phương sai (variance) tính trên cửa sổ trượt cực nhỏ là đủ để phân biệt cú ngã, mở ra hướng tối giản hóa thuật toán cho hệ thống nhúng. Điểm hạn chế là dataset chỉ có một tình nguyện viên nam trẻ tuổi, chưa phản ánh được động lực học té ngã của người cao tuổi.

\end{enumerate}

\subsubsection{Phương pháp học máy}
\label{subsubsection:2.1.2.b}

Các phương pháp học máy (machine learning) cải thiện độ chính xác phát hiện bằng cách học các mẫu đặc trưng phức tạp từ dữ liệu huấn luyện thay vì dựa trên ngưỡng cứng. Pipeline điển hình gồm tiền xử lý tín hiệu, trích xuất đặc trưng thủ công và huấn luyện một mô hình phân loại. Tuy độ chính xác cao hơn nhóm threshold-based, chi phí tính toán trong pha suy luận vẫn đủ nhỏ để chạy trên thiết bị nhúng.

\begin{enumerate}[label=\arabic*)]

\item Nazar và Jalal \cite{nazar2025csohm} đề xuất hệ thống phát hiện té ngã kết hợp thuật toán tối ưu bầy đàn Cuckoo (Cuckoo Search Optimization --- CSO) để lựa chọn đặc trưng với mô hình Markov ẩn dạng Gaussian (Hidden Markov Model --- HMM) để phân loại chuỗi thời gian. Dữ liệu IMU được khử nhiễu bằng bộ lọc Kalman trước khi phân đoạn qua cửa sổ trượt 200 mẫu; các đặc trưng entropy Shannon, độ lệch chuẩn, năng lượng phổ, chiều dài tích lũy và biên độ tín hiệu được tối ưu bằng CSO để loại bỏ dư thừa. Mỗi hoạt động được đại diện bởi một HMM ba trạng thái ẩn, giúp nắm bắt tốt sự phụ thuộc tuần tự của tín hiệu chuyển động và phân biệt các kiểu ngã có động lực học gần giống nhau (trượt chân, vấp ngã, ngất xỉu). Kiểm thử trên tập KFall đạt độ chính xác 93,7\% cho 15 lớp hoạt động, vượt trội SVM (84,89\%) và Random Forest (89,37\%). Hạn chế là pha tối ưu hóa CSO và huấn luyện HMM bằng thuật toán Baum-Welch yêu cầu tài nguyên tính toán rất lớn, bắt buộc phải chạy offline, không phù hợp triển khai trực tiếp trên vi điều khiển cạnh.

\item Aqillah và cộng sự \cite{aqillah2024comparative} thực hiện nghiên cứu so sánh bốn mô hình học máy trên cùng một hệ thống phát hiện té ngã thời gian thực dùng cảm biến IMU: Máy véc-tơ hỗ trợ (SVM), Rừng ngẫu nhiên (RF), Mạng nơ-ron tích chập (CNN) và Mạng nơ-ron nhân tạo (ANN). Dữ liệu gia tốc, vận tốc góc và góc định hướng được thu thập ở 100 Hz từ các hoạt động giả lập gồm đi bộ, đứng yên và ngã theo nhiều hướng. ANN đạt hiệu quả tốt nhất với độ chính xác 98\% và độ trễ suy luận chỉ 0,1 giây, tiếp theo là CNN (97\%), RF (96\%) và SVM (90\%). Hệ thống được thiết kế để kích hoạt phản ứng bảo vệ tức thì (bung túi khí vest) khi phát hiện ngã, nhấn mạnh yêu cầu khắt khe về độ trễ thấp trong ứng dụng y tế khẩn cấp. Nghiên cứu này là một trong những cơ sở định lượng để đồ án lựa chọn kiến trúc ResNet 1D thay vì SVM/RF; điểm hạn chế là dataset nhỏ với chỉ một người thực hiện giả lập, chưa có dữ liệu người cao tuổi.

\item Park và cộng sự \cite{park2025fallrisk} nghiên cứu chiến lược tối ưu hóa vị trí đặt cảm biến IMU và thuật toán phân loại để tối đa hóa độ chính xác đánh giá nguy cơ té ngã thông qua phân tích dáng đi. Nghiên cứu thu thập dữ liệu động học từ 93 người cao tuổi sống trong cộng đồng (tuổi trung bình 77,8) sử dụng 10 cảm biến IMU đặt tại các vị trí khác nhau trên cơ thể trong bài kiểm tra đi bộ 10 mét, với đánh giá nguy cơ ngã qua công cụ PPA rút gọn. Kết quả cho thấy vị trí đặt cảm biến có ảnh hưởng có ý nghĩa thống kê đến độ chính xác phân loại ($p < 0{,}001$), trong khi sự khác biệt giữa sáu thuật toán (SVM, DT, RF, KNN, XGBoost, LightGBM) là không đáng kể. Hiệu quả tốt nhất đạt được bởi SVM dùng cảm biến chân dưới bên trái với độ chính xác 90,1\%, độ nhạy 95,7\% và độ đặc hiệu 84,1\%. Phát hiện quan trọng nhất của nghiên cứu --- vùng thắt lưng và chân dưới cho tín hiệu dáng đi chất lượng cao nhất với người cao tuổi --- củng cố trực tiếp quyết định đặt thiết bị đeo tại thắt lưng trong đồ án.

\end{enumerate}

\subsubsection{Phương pháp học sâu}
\label{subsubsection:2.1.2.c}

Học sâu (deep learning) cho phép học đặc trưng tự động từ dữ liệu thô mà không cần bước trích xuất đặc trưng thủ công, đạt độ chính xác vượt trội trong các bài toán phân loại tín hiệu chuỗi thời gian phức tạp. Hạn chế chính là yêu cầu bộ nhớ và tính toán lớn, đặt ra thách thức triển khai lên vi điều khiển nhúng.

\begin{enumerate}[label=\arabic*)]

\item Cai và cộng sự \cite{cai2026stanet} xây dựng tập dữ liệu quy mô lớn FallTL (28,4 giờ dữ liệu, 45 người tham gia, 8 vị trí đặt cảm biến) và phát triển mạng STA-Net (Spatial-Temporal Attention Network) kiến trúc nhánh kép để phát hiện té ngã trước thời điểm va chạm (Pre-Impact Fall Detection --- PIFD). Nhánh không gian dùng tích chập 2D đa tỷ lệ kết hợp khối chú ý SE (Squeeze-and-Excitation) để học tương quan giữa các trục cảm biến, trong khi nhánh thời gian dùng Bi-GRU kết hợp Criss-Cross Attention để nắm bắt phụ thuộc dài hạn theo thời gian. Mô hình huấn luyện với hàm mất mát Focal Loss nhằm xử lý mất cân bằng lớp nghiêm trọng giữa mẫu hoạt động bình thường và mẫu té ngã. Kết quả đạt khoảng 94\% (PIFD trên FallTL) và 96--99\% (phát hiện sau va chạm), với thời gian cảnh báo sớm trung bình 456,69 ms trước khi chạm đất --- đủ để kích hoạt thiết bị bảo vệ chủ động. Hạn chế chính là kích thước mô hình 22,4 MB vượt quá giới hạn Flash/RAM của đa số vi điều khiển, đòi hỏi lượng tử hóa INT8 hoặc chạy trên Edge Gateway để triển khai thực tế.

\item Mekruksavanich và cộng sự \cite{mekruksavanich2024resnet} đề xuất kiến trúc mạng phần dư một chiều ResNeXt tùy chỉnh (1D ResNeXt) để phân tích tín hiệu IMU từ gia tốc kế và con quay hồi chuyển ba trục. Kiến trúc phần dư (residual) giải quyết vấn đề tiêu biến gradient trong quá trình huấn luyện mạng sâu thông qua các kết nối tắt (skip connection), cho phép học các đặc trưng phân cấp từ tín hiệu chuyển động. Hệ thống được kiểm thử trên tập PIF gồm 12 cá thể đa dạng về độ tuổi và thể trạng, đạt độ chính xác 91,52\% chỉ với gia tốc kế, 83,80\% với con quay hồi chuyển, và 92,27\% khi kết hợp cả hai --- vượt trội các baseline CNN, LSTM, BiLSTM, GRU và BiGRU. Đây là nền tảng trực tiếp cho quyết định lựa chọn kiến trúc ResNet 1D trong đồ án để chạy suy luận on-device trên ESP32-S3; điểm cần lưu ý là tập PIF còn nhỏ (12 người) và chỉ gồm dữ liệu ngã giả lập từ người trẻ.

\item Gaud và cộng sự \cite{gaud2025fibils} đề xuất mô hình FIBiLS kết hợp mạng nơ-ron tích chập 1D (1D-CNN) và mạng LSTM hai chiều (BiLSTM) để phát hiện sớm té ngã trong quá trình thực hiện hoạt động hằng ngày. Hệ thống sử dụng hai cảm biến IMU đeo trên người và huấn luyện trên tập dữ liệu công khai SisFall --- bộ dữ liệu gồm 19 loại hoạt động hằng ngày và 15 kiểu té ngã từ 23 người trẻ và 14 người cao tuổi trên 62 tuổi, thu thập bằng thiết bị tự chế với gia tốc kế và con quay hồi chuyển \cite{sucerquia2017sisfall}. FIBiLS đạt độ chính xác 99,68\% với thời gian suy luận nhanh và được triển khai thực tế lên vi điều khiển NodeMCU để kiểm chứng khả năng tính toán biên. Sự kết hợp giữa đặc trưng cục bộ (CNN) và phụ thuộc tuần tự (BiLSTM) cho phép mô hình nắm bắt đồng thời hình dạng tín hiệu và tiến trình chuyển động theo thời gian; hạn chế là tập SisFall thiên lệch về dữ liệu người trẻ, ảnh hưởng đến hiệu quả nhận diện với người cao tuổi thực sự.

\end{enumerate}

\subsubsection{Nhận xét tổng hợp}
\label{subsubsection:2.1.2.d}

Bảng \ref{tab:survey_compare} tổng hợp và so sánh các công trình đã khảo sát theo các tiêu chí phù hợp với yêu cầu của đồ án.

\begin{table}[H]
    \centering
    \caption{Bảng so sánh các công trình nghiên cứu phát hiện té ngã}
    \label{tab:survey_compare}
    \small
    \begin{tabular}{|p{2.7cm}|p{2.6cm}|p{2.5cm}|p{2.0cm}|p{2.0cm}|p{2.9cm}|}
        \hline
        \textbf{Tác giả \& Năm} & \textbf{Phương pháp} & \textbf{Phần cứng / Dataset} & \textbf{Độ chính xác} & \textbf{On-device?} & \textbf{Hướng áp dụng} \\
        \hline
        Alsaadi \textit{et al.}, 2024 & Logic tree + SVR thích ứng & 4 IMU 9 trục, 20 người & 90,8--92,1\% & Có (MCU) & Pipeline tiền xử lý, RMS triggering \\
        \hline
        Tseng \textit{et al.}, 2025 & FSM + GPS + NB-IoT & IMU+GPS, 15 người, 6.750 mẫu & 99,7\% & Có & FSM event-driven, firmware ESP32 \\
        \hline
        Booranawong \textit{et al.}, 2025 & KNN + EWMA (cầu thang) & GY-521 + XBee3, 1 người & 88\% TB; 100\% ngã & Có (MCU) & EWMA filter, variance feature \\
        \hline
        Nazar \& Jalal, 2025 & CSO + HMM & KFall, IMU 6 trục & 93,7\% & Không & Kalman filter, lựa chọn đặc trưng \\
        \hline
        Aqillah \textit{et al.}, 2024 & SVM / RF / CNN / ANN & IMU 6 trục, 100 Hz & ANN: 98\%; CNN: 97\% & Hạn chế & Cơ sở chọn ResNet thay SVM \\
        \hline
        Park \textit{et al.}, 2025 & SVM, phân tích vị trí cảm biến & 10 IMU, 93 người cao tuổi & 90,1\% & Có & Khẳng định đặt cảm biến thắt lưng \\
        \hline
        Cai \textit{et al.}, 2026 & STA-Net (Dual-branch, PIFD) & FallTL + KFall, 45 người & PIFD: 94\%; FD: 99\% & Không (22 MB) & Focal Loss, sliding window 1 s \\
        \hline
        Mekruksavanich \textit{et al.}, 2024 & 1D ResNeXt & PIF, 12 người & 92,27\% & Cần tối ưu & \textbf{Kiến trúc ResNet 1D đồ án} \\
        \hline
        Gaud \textit{et al.}, 2025 & FIBiLS (CNN + BiLSTM) & SisFall, 37 người, 2 IMU & 99,68\% & Có (NodeMCU) & Edge deployment, SisFall dataset \\
        \hline
    \end{tabular}
\end{table}

Qua khảo sát, có thể rút ra một số nhận xét định hướng quan trọng cho đồ án. Về \textbf{phần cứng}, vị trí đặt cảm biến tại thắt lưng được xác nhận bởi nhiều nghiên cứu \cite{booranawong2025stairs, cai2026stanet, park2025fallrisk} là lựa chọn tối ưu về chất lượng tín hiệu lẫn sự thuận tiện cho người dùng cao tuổi. Tổ hợp gia tốc kế và con quay hồi chuyển (6 trục) là đủ cho bài toán phát hiện ngã, không cần từ kế \cite{cai2026stanet}. Về \textbf{thuật toán}, phương pháp học sâu với kiến trúc mạng phần dư (ResNet) một chiều đạt cân bằng tốt nhất giữa độ chính xác và khả năng triển khai trên vi điều khiển sau khi lượng tử hóa, so với các mô hình lớn như STA-Net hay FIBiLS BiLSTM. Về \textbf{dataset}, các tập dữ liệu hiện có chủ yếu được thu thập từ người trẻ giả lập, là điểm yếu chung của toàn ngành, đặt ra thách thức về tính khái quát hóa cho người cao tuổi. Dựa trên phân tích trên, đồ án lựa chọn xây dựng mô hình phát hiện té ngã theo kiến trúc \textbf{ResNet 1D} triển khai trên ESP32-S3, kết hợp với kiến trúc truyền thông hướng sự kiện qua MQTT 4G LTE để đáp ứng đồng thời yêu cầu về độ chính xác, độ trễ thấp và khả năng triển khai thực tế trong môi trường cơ sở y tế.

\section{Tổng quan chức năng}
\label{section:2.2}
Phần này mô tả tổng quan các chức năng cốt lõi của hệ thống, dựa trên định hướng gồm một vai trò người dùng (Nhân viên y tế) và một tác nhân hệ thống (Thiết bị giám sát).

\subsection{Biểu đồ use case tổng quát}
\label{subsection:2.2.1}
Biểu đồ use case tổng quát mô tả các tương tác chính giữa con người và hệ thống. Các tác nhân (Actors) tham gia bao gồm:
\begin{itemize}
    \item \textbf{Nhân viên y tế (Medical Staff):} Người sử dụng hệ thống qua giao diện Dashboard để giám sát thời gian thực, theo dõi hồ sơ bệnh nhân, thiết lập/quản lý thiết bị và xử lý các sự cố té ngã khẩn cấp.
    \item \textbf{Thiết bị giám sát (ESP32 Node):} Tác nhân hệ thống (System Actor), liên tục thu thập và gửi dữ liệu viễn trắc (telemetry) thông qua giao thức PPPoS trên mạng 4G LTE và kích hoạt các cảnh báo (alerts) khẩn cấp.
\end{itemize}

Các ca sử dụng (Use case) cốt lõi bao gồm: Đăng nhập \& Xác thực, Quản lý thiết bị, Quản lý bệnh nhân (Wearer), Gán/Thu hồi thiết bị, Giám sát thời gian thực, Xử lý cảnh báo té ngã, và Xem thống kê vận động.

\subsection{Biểu đồ use case phân rã "Giám sát thời gian thực"}
\label{subsection:2.2.2}
Use case "Giám sát thời gian thực" yêu cầu hệ thống cung cấp cho nhân viên y tế một giao diện Dashboard tổng quan. Thông qua giao diện này, nhân viên y tế có thể theo dõi danh sách toàn bộ các thiết bị đang hoạt động, xem trạng thái kết nối mạng, mức pin hiện tại của từng thiết bị. Đặc biệt, chức năng này ứng dụng cơ chế \textbf{Đồng bộ cảnh báo lai (Hybrid Alert Sync)}, liên tục lắng nghe các gói tin từ MQTT Broker qua WebSocket để ngay lập tức hiển thị cảnh báo (kèm âm thanh và viền đỏ nhấp nháy) khi có sự cố té ngã xảy ra với độ trễ chưa tới 1 giây, đồng thời tự động cập nhật dữ liệu từ API. Biểu đồ phân rã cho nhóm chức năng này tập trung vào luồng hiển thị dữ liệu và cơ chế cảnh báo.

% \begin{figure}[H]
%     \centering
%     % \includegraphics[width=0.8\textwidth]{Hinhve/use_case_giam_sat.png}
%     \caption{Biểu đồ use case phân rã "Giám sát thời gian thực"}
%     \label{fig:use_case_giam_sat}
% \end{figure}

\subsection{Biểu đồ use case phân rã "Xử lý cảnh báo té ngã"}
\label{subsection:2.2.3}
Use case "Xử lý cảnh báo té ngã" mô tả quy trình tiếp nhận và xử lý sự cố. Khi màn hình Dashboard hiển thị cảnh báo té ngã, nhân viên y tế sẽ kiểm tra và đến hiện trường để hỗ trợ người bệnh. Sau khi hoàn thành sơ cứu hoặc xác minh tình trạng an toàn, nhân viên y tế sẽ tương tác với hệ thống bằng cách nhấn nút "Đã xử lý" (Resolve). Chức năng này cập nhật trạng thái của cảnh báo trong cơ sở dữ liệu để đảm bảo không có sự cố nào bị bỏ sót, đồng thời lưu trữ lịch sử để phục vụ báo cáo sau này.

% \begin{figure}[H]
%     \centering
%     % \includegraphics[width=0.8\textwidth]{Hinhve/use_case_xu_ly.png}
%     \caption{Biểu đồ use case phân rã "Xử lý cảnh báo té ngã"}
%     \label{fig:use_case_xu_ly}
% \end{figure}

\subsection{Biểu đồ use case phân rã "Gán thiết bị cho bệnh nhân"}
\label{subsection:2.2.4}
Use case "Gán thiết bị cho bệnh nhân" phục vụ mục đích thiết lập liên kết vật lý giữa một thiết bị ESP32 và một hồ sơ bệnh nhân. Nhân viên y tế sẽ chọn một thiết bị chưa được sử dụng và chọn một bệnh nhân từ danh sách. Một yêu cầu cực kỳ quan trọng ở bước này là việc nhập chiều cao của bệnh nhân, vì hệ thống cần thông số này để ước tính sải bước chân và tính toán quãng đường di chuyển dựa trên dữ liệu gia tốc. Việc gán thiết bị cũng tự động kích hoạt quá trình thu thập và lưu trữ viễn trắc cho bệnh nhân đó.

% \begin{figure}[H]
%     \centering
%     % \includegraphics[width=0.8\textwidth]{Hinhve/use_case_gan_thiet_bi.png}
%     \caption{Biểu đồ use case phân rã "Gán thiết bị cho bệnh nhân"}
%     \label{fig:use_case_gan_thiet_bi}
% \end{figure}

\subsection{Biểu đồ use case phân rã "Xem thống kê vận động"}
\label{subsection:2.2.5}
Use case "Xem thống kê vận động" giúp nhân viên y tế và quản trị viên tra cứu các số liệu lịch sử. Người dùng có thể chọn một hồ sơ bệnh nhân cụ thể và thiết lập mốc thời gian cần tra cứu (ví dụ: trong tuần qua). Hệ thống sẽ truy xuất dữ liệu chuỗi thời gian (time-series) từ InfluxDB và vẽ các biểu đồ trực quan như: biểu đồ số bước chân đi bộ, biểu đồ số bước chạy, và biểu đồ khoảng cách di chuyển. Các biểu đồ này cung cấp thông tin y khoa quan trọng để đánh giá mức độ phục hồi hoặc suy giảm chức năng vận động của người bệnh.

% \begin{figure}[H]
%     \centering
%     % \includegraphics[width=0.8\textwidth]{Hinhve/use_case_thong_ke.png}
%     \caption{Biểu đồ use case phân rã "Xem thống kê vận động"}
%     \label{fig:use_case_thong_ke}
% \end{figure}

\subsection{Quy trình nghiệp vụ (Activity Diagram)}
\label{subsection:2.2.6}
Quy trình nghiệp vụ quan trọng nhất của hệ thống là \textbf{Luồng xử lý sự cố té ngã khẩn cấp}. Quy trình diễn ra như sau:
\begin{enumerate}
    \item Thiết bị đeo (ESP32-S3) phát hiện té ngã thông qua mô hình học sâu 1D ResNet.
    \item Thiết bị ngay lập tức gửi gói tin cảnh báo qua giao thức MQTT (kết nối 4G LTE bằng PPPoS) đến MQTT Broker.
    \item Giao diện Web (Dashboard) đang đăng ký (subscribe) luồng MQTT sẽ lập tức hiển thị cảnh báo đỏ toàn màn hình kèm theo âm thanh báo động.
    \item Nhân viên y tế xác nhận cảnh báo trên màn hình và di chuyển đến vị trí bệnh nhân để sơ cứu.
    \item Sau khi xử lý xong sự cố, nhân viên y tế thao tác chọn "Đã xử lý" (Resolve) trên hệ thống.
    \item Hệ thống gọi API cập nhật trạng thái cảnh báo vào cơ sở dữ liệu PostgreSQL, kết thúc quy trình.
\end{enumerate}

\section{Đặc tả chức năng}
\label{section:2.3}
Dưới đây là đặc tả chi tiết của 4 Use Case quan trọng nhất, phản ánh nghiệp vụ cốt lõi của hệ thống giám sát.

\subsection{Đặc tả Use case: Giám sát thời gian thực}
\begin{table}[H]
    \centering
    \caption{Đặc tả Use case: Giám sát thời gian thực}
    \label{tab:uc_giam_sat}
    \begin{tabular}{|p{4cm}|p{10.5cm}|}
        \hline
        \textbf{Tên Use case} & \textbf{Giám sát thời gian thực} \\ \hline
        \textbf{Tác nhân} & Nhân viên y tế \\ \hline
        \textbf{Mô tả} & Cung cấp giao diện lưới (grid) hiển thị trạng thái của toàn bộ thiết bị đang hoạt động (Pin, kết nối mạng) và tiếp nhận cảnh báo khẩn cấp. \\ \hline
        \textbf{Tiền điều kiện} & Nhân viên đã đăng nhập và có thiết bị đang online. \\ \hline
        \textbf{Luồng sự kiện chính} &
        1. Nhân viên y tế truy cập vào trang chủ (Dashboard). \newline
        2. Hệ thống (Frontend) thiết lập kết nối WebSocket tới MQTT Broker. \newline
        3. Hệ thống hiển thị danh sách thiết bị kèm trạng thái kết nối. \newline
        4. Khi có gói tin cảnh báo té ngã từ Broker, Frontend lập tức bung màn hình cảnh báo (Overlay) và phát âm thanh. \\ \hline
        \textbf{Hậu điều kiện} & Cảnh báo được hiển thị thành công cho đến khi nhân viên xác nhận. \\ \hline
    \end{tabular}
\end{table}

\subsection{Đặc tả Use case: Xử lý cảnh báo té ngã}
\begin{table}[H]
    \centering
    \caption{Đặc tả Use case: Xử lý cảnh báo té ngã (Resolve Alert)}
    \label{tab:uc_xu_ly_canh_bao}
    \begin{tabular}{|p{4cm}|p{10.5cm}|}
        \hline
        \textbf{Tên Use case} & \textbf{Xử lý cảnh báo té ngã} \\ \hline
        \textbf{Tác nhân} & Nhân viên y tế \\ \hline
        \textbf{Mô tả} & Nhân viên đánh dấu một cảnh báo té ngã là "Đã xử lý" sau khi đã hỗ trợ bệnh nhân. \\ \hline
        \textbf{Tiền điều kiện} & Tồn tại ít nhất một cảnh báo té ngã chưa được xử lý. \\ \hline
        \textbf{Luồng sự kiện chính} &
        1. Nhân viên y tế bấm vào nút "Xác nhận đã xử lý" trên cảnh báo tương ứng. \newline
        2. Frontend gửi yêu cầu PATCH kèm Access Token tới Backend FastAPI. \newline
        3. Backend cập nhật trạng thái \texttt{is\_resolved = True} trong cơ sở dữ liệu PostgreSQL. \newline
        4. Backend trả về phản hồi thành công, Frontend ẩn cảnh báo khỏi danh sách chờ. \\ \hline
        \textbf{Hậu điều kiện} & Cảnh báo được lưu trữ vào lịch sử với trạng thái đã xử lý. \\ \hline
    \end{tabular}
\end{table}

\subsection{Đặc tả Use case: Gán thiết bị cho bệnh nhân}
\begin{table}[H]
    \centering
    \caption{Đặc tả Use case: Gán thiết bị cho bệnh nhân (Assign Device)}
    \label{tab:uc_gan_thiet_bi}
    \begin{tabular}{|p{4cm}|p{10.5cm}|}
        \hline
        \textbf{Tên Use case} & \textbf{Gán thiết bị cho bệnh nhân} \\ \hline
        \textbf{Tác nhân} & Nhân viên y tế \\ \hline
        \textbf{Mô tả} & Liên kết một thiết bị giám sát vật lý với hồ sơ của một bệnh nhân cụ thể. \\ \hline
        \textbf{Tiền điều kiện} & Thiết bị đang ở trạng thái trống (chưa gán) và hồ sơ bệnh nhân đã được tạo. \\ \hline
        \textbf{Luồng sự kiện chính} &
        1. Nhân viên y tế chọn thiết bị từ danh sách và nhấn "Gán bệnh nhân". \newline
        2. Hệ thống hiển thị danh sách bệnh nhân chưa có thiết bị. \newline
        3. Nhân viên chọn bệnh nhân tương ứng và bấm lưu. \newline
        4. Hệ thống Backend kiểm tra ràng buộc 1-1 (unique constraint) và cập nhật \texttt{wearer\_id} cho thiết bị. \\ \hline
        \textbf{Hậu điều kiện} & Dữ liệu viễn trắc của thiết bị từ thời điểm này sẽ được tính toán dựa trên chiều cao của bệnh nhân được gán. \\ \hline
    \end{tabular}
\end{table}

\subsection{Đặc tả Use case: Xem thống kê vận động}
\begin{table}[H]
    \centering
    \caption{Đặc tả Use case: Xem thống kê vận động}
    \label{tab:uc_thong_ke}
    \begin{tabular}{|p{4cm}|p{10.5cm}|}
        \hline
        \textbf{Tên Use case} & \textbf{Xem thống kê vận động} \\ \hline
        \textbf{Tác nhân} & Nhân viên y tế \\ \hline
        \textbf{Mô tả} & Tra cứu và xem biểu đồ thống kê mức độ vận động (số bước chân, quãng đường) theo tuần/tháng. \\ \hline
        \textbf{Tiền điều kiện} & Bệnh nhân đã đeo thiết bị và sinh ra dữ liệu viễn trắc. \\ \hline
        \textbf{Luồng sự kiện chính} &
        1. Người dùng chọn khoảng thời gian cần tra cứu trong tab Thống kê. \newline
        2. Backend tiến hành truy vấn (Query) dữ liệu tổng hợp từ InfluxDB. \newline
        3. Dữ liệu được trả về và Frontend sử dụng thư viện Recharts để vẽ biểu đồ trực quan. \\ \hline
        \textbf{Hậu điều kiện} & Biểu đồ được hiển thị chính xác. \\ \hline
    \end{tabular}
\end{table}

\section{Yêu cầu phi chức năng}
\label{section:2.4}
Hệ thống giám sát khẩn cấp đòi hỏi các tiêu chuẩn khắt khe về kỹ thuật và độ tin cậy:
\begin{itemize}
    \item \textbf{Độ trễ thấp (Low Latency):} Thời gian từ lúc thiết bị phát hiện té ngã đến khi chuông báo động vang lên trên màn hình Web của y tá phải dưới 1 giây.
    \item \textbf{Độ tin cậy truyền thông:} Sử dụng ngăn xếp giao thức mạng qua giao tiếp PPPoS (Point-to-Point Protocol over Serial) trên module SIM 4G kết hợp cơ chế QoS của MQTT để đảm bảo gói tin cảnh báo luôn được gửi đi ngay cả trong điều kiện sóng yếu.
    \item \textbf{Khả năng mở rộng (Scalability):} Kiến trúc Dual-Database giúp hệ thống xử lý đồng thời hàng nghìn luồng dữ liệu viễn trắc (InfluxDB) mà không gây tắc nghẽn cho các truy vấn quản lý (PostgreSQL). Hỗ trợ kiến trúc đa khách thuê (Multi-tenancy) cho phép quản lý đồng thời nhiều viện dưỡng lão trên cùng một máy chủ vật lý.
    \item \textbf{Tính bảo mật (Security):} Mọi API giao tiếp từ Frontend đều phải được xác thực thông qua JSON Web Token (JWT), kết hợp phân quyền quản trị đa tổ chức, chống rò rỉ thông tin cá nhân của bệnh nhân.
\end{itemize}

\end{document}
